% amsmath is already loaded by the theme. So is a sans maths font, so
% symbols match the body weight instead of sitting in a serif face.

% --- Inline and display --------------------------------------------
The error level $\lambda$ controls it.

\[ q_j = \frac{e^{\lambda \mathrm{EU}_j(x)}}{\sum_i e^{\lambda \mathrm{EU}_i(x)}} \]

% Display maths is UNNUMBERED by default. A talk has no cross-references.

% --- When you must number one --------------------------------------
\begin{equation}
  \mathcal{L} = \mathbb{E}_{\pi}\left[ \sum_t \gamma^t r_t \right]
  \label{eq:objective}
\end{equation}
Referred to later as \eqref{eq:objective}.

% --- Several lines, aligned on = -----------------------------------
\begin{align*}
  \mathcal{L} &= \mathbb{E}_{\pi}\left[ \sum_t \gamma^t r_t \right] \\
              &= \sum_s d^{\pi}(s) \, v^{\pi}(s)
\end{align*}
% align* is unnumbered; align numbers every line.

% --- Cases and matrices --------------------------------------------
\[ r(s) = \begin{cases}
     +1 & \text{host recovered} \\
     -1 & \text{host compromised} \\
      0 & \text{otherwise}
   \end{cases} \]

\[ A = \begin{bmatrix} a & b \\ c & d \end{bmatrix} \]

% --- Pointing at part of an equation -------------------------------
% Colour the piece you are talking about:
\[ q_j = \frac{e^{\tint{accent}{\lambda} \mathrm{EU}_j(x)}}
              {\sum_i e^{\lambda \mathrm{EU}_i(x)}} \]

% Or reveal it, then emphasise it, across two pages:
\begin{frame}{Building the objective}
  \[ \mathcal{L} = \emphat{2}{\mathbb{E}_{\pi}}\left[ \sum_t \gamma^t r_t \right] \]
  \showfrom{2}{The expectation is over the policy, not the data.}
\end{frame}

% --- Statements ----------------------------------------------------
\begin{theorem}
  The equilibrium exists for every finite game.
\end{theorem}

\begin{definition}
  A \term{world model} predicts the next observation from the last.
\end{definition}

% Also available: lemma, corollary, proof, example, block.
% Give any of them a name: \begin{theorem}[Nash, 1951]

% --- Two habits worth keeping --------------------------------------
% \mathrm{} for words inside maths: \mathrm{EU}, not EU
% \dfrac for a full-size fraction inside inline maths
